\begingroup
\setlength{\emergencystretch}{2em}
\capituloTesis{Recomendaciones}

\seccionTesis{Recomendación vinculada al objetivo general}

Se recomienda mantener diferenciadas la evacuación longitudinal, la abertura transversal, la condición superficial y la rugosidad en los futuros registros de conservación vial del corredor, conservando para cada unidad estructura--sector su identificación, progresivas, coordenadas y respaldo instrumental. Dado que solo una de las cuatro asociaciones cumplió el criterio estadístico familiar, la evaluación conjunta debe presentar las relaciones por separado en lugar de sintetizarlas en un índice o decisión global no justificados. Para investigaciones posteriores que busquen contrastar una relación general entre las dos variables, corresponde definir previamente un modelo multivariado o un criterio de integración y comprobar sus supuestos con una muestra apropiada.

\seccionTesis{Recomendación vinculada al primer objetivo específico}

Se recomienda priorizar la revisión conjunta de los recorridos longitudinales y del estado superficial del pavimento en las unidades donde el registro indique baja continuidad de evacuación, sin interpretar esa coincidencia como demostración de que la discontinuidad haya causado los deterioros. La inspección debería conservar por margen la longitud evaluada, la longitud continua, la ubicación de interrupciones y la evidencia fotográfica, y contrastar esas observaciones con los tipos, severidades y extensiones de falla utilizados para calcular el PCI. En evaluaciones sucesivas conviene registrar además tránsito, intervenciones de mantenimiento y características del pavimento, a fin de examinar factores que pueden influir en la asociación identificada.

\seccionTesis{Recomendación vinculada al segundo objetivo específico}

Se recomienda repetir la evaluación de CL e IRI en campañas comparables y con mayor cobertura de unidades, debido a que la correlación inversa observada no conservó significancia después del control familiar por Holm. La toma de datos debería mantener perfiles longitudinales identificables por huella, recorrido, progresiva y fecha, y documentar cambios en la continuidad de evacuación y en las condiciones de la calzada; de este modo podrá evaluarse si la dirección y magnitud de CL--IRI se reproducen bajo distintos momentos y sectores. La planificación de nuevas pruebas deberá establecer con anticipación el tamaño muestral, la dependencia espacial y el control de multiplicidad.

\seccionTesis{Recomendación vinculada al tercer objetivo específico}

Se recomienda conservar la inspección independiente de la entrada y la salida de cada estructura transversal, registrando las áreas interiores, las porciones ocupadas y las áreas libres de ambos extremos, junto con su ubicación y evidencia fotográfica. Como la asociación CT--PCI no alcanzó significancia, no corresponde utilizar únicamente el porcentaje de abertura libre para inferir la condición superficial del pavimento ni para establecer una prioridad de intervención. Cuando el propósito sea estimar el funcionamiento hidráulico de los pases, deberá complementarse la geometría registrada con información de pendientes, caudales de diseño, control de entrada y salida y otros parámetros exigidos por el análisis hidráulico.

\seccionTesis{Recomendación vinculada al cuarto objetivo específico}

Se recomienda mantener evaluaciones independientes de CT e IRI y ampliar la observación a sectores y momentos adicionales antes de formular decisiones sustentadas en su asociación, puesto que el coeficiente fue pequeño, no significativo y cambió de signo al excluir determinadas unidades. El registro de rugosidad debería preservar los perfiles de las dos huellas y sus comprobaciones de repetición, mientras que la inspección transversal debería documentar el estado de cada extremo, la condición de descarga y los cambios de obstrucción. Si una investigación posterior pretende explicar la rugosidad a partir del drenaje transversal, será necesario incorporar otras características del pavimento y de la vía dentro de un diseño analítico previamente especificado.

\endgroup
